Télécom Paris, a founding member of Institut Polytechnique de Paris and a part of IMT (Institut Mines-Télécom), is ranked as one of the top 5 engineering schools in France. Recognized for its invaluable corporate connections, Telecom Paris is an institution with a strong sense of community and international representation. This public higher-education establishment guarantees excellent employability in all sectors and is the leading engineering school for the entire digital spectrum (from hardware to software technology).

With its high-level instruction and innovative pedagogy, Télécom Paris is at the heart of a unique innovation ecosystem which is based on collaborative classroom engagement, an emphasis on project-based training, and multidisciplinary research. Our scientific faculty are affiliated with two key research laboratories; first, the LTCI laboratory, which is recognized by the HCERES as a leading facility in the field of digital sciences and has a remarkable international outreach. And second, the i3 laboratory, Interdisciplinary Institute of Innovation (I3 - UMR 9217 of the CNRS), which carries out multidisciplinary research focused on innovation and in collaboration with the École Polytechnique and Mines ParisTech. 

Based in Palaiseau, France at the heart of the Institut Polytechnique campus alongside the École Polytechnique, ENSTA, Télécom Sud Paris and ENSAE, Télécom Paris also has a central-Paris incubator at the epicenter of the French start-up ecosystem. 

Télécom Paris is a living laboratory for all the major technological and societal challenges including Artificial Intelligence, Quantum Computing, IoT, Cybersecurity, large-scale digital equipment (Cloud Computing), 5G/6G, and Green IT.